\scanpage{25}
Pour le voir, \deleted{\unsure{c'est comme aussitôt}} \inserted{comme} c'est évidemment injectif, on est ramené à prouver que c'est surjectif. Soit $g':S'\to X$, prouvons que ça se factorise par $S$. \deleted{Or \unsure{$g$ définit une}} \inserted{Comme $S'\to S$ est un homéomorphisme, on est ramené trivialement au cas où $S$, donc $S'$, local, de pt} \deleted{section de $X'=X\times_S S'$ sur $S'$, donc un point $x'$ de $X'$, \inserted{or $X'/S'$ est \deleted{\unsure{form}} \unsure{form\textsuperscript{t} net}} (comme $X'\to X$ est un homéomorphisme, on conclut que $g(S')\inserted{=U}$ est une partie \unsure{et} à la fois ouverte de} \inserted{pt fermé $s,s'$, et à prouver que (si $x=g'(s')$) $\underline O_{X,x}\to\mathcal O_{S',s'}$ se factorise par $\mathcal O_{S,s}$.}
\diagram{\begin{tikzcd}[ampersand replacement=\&,column sep=small,row sep=small]
\& \underline O_{X,x}\arrow[dl,dashed]\\
\underline O_{S',s'} \& \mathcal O_{S,s}\arrow[l,hook]\arrow[u,hook]
\end{tikzcd}}
\deleted{\unsure{fermée de} $X$, et \unsure{quitte à remplacer} $X$ par $U$, on peut} \inserted{remplaçant $X$ par un voisinage affine de $x$, on peut supposer $X$ affine, donc} \deleted{supposer que $X\to S$ est un homéomorphisme, donc} séparé, et \unsure{par suite|sous} \deleted{quitte à remplacer $X$ par \unsure{$g(S)$} $g'(S')$ \unsure{soit}} $g'$ est \deleted{\unsure{partie}}, \deleted{une immer. fermée, Mais dans ce cas on est} une \unsure{fermée}. Remplaçant $X$ par l'image \deleted{fermée} \inserted{schématique} de $g'$, on peut supposer \deleted{Mais si $\mathcal A=f_\ast(\mathcal O_{S'})$, \unsure{on a} $\mathcal J\mathcal A=\unsure{0}$ et comme $\mathcal O_S\to\mathcal A$ est injectif, on en conclut $\mathcal J=0$, cqfd.} \inserted{que $X$} est \deleted{\unsure{donc}} \deleted{\unsure{alors sur $S$}} local, \deleted{\unsure{donc}} de pt fermé $x$, $\underline O_{X,x}\hookrightarrow\mathcal O_{S',s'}$. Alors $\underline O_{X,x}$ est \deleted{donc} entier sur $\mathcal O_{S,\unsure{s|s'}}$, et \unsure{$=$|même} fini sur $\mathcal O_{S,s}$ si $f$ fini et $S$ loc. noeth. $\ldots$

\margin{\rule{.6pt}{2.6\baselineskip}}\uline{Corollaire 2} Soit $f:\inserted{S'\to S}$ un épimorphisme propre \inserted{\deleted{\unsure{de prés. finie|loc\textsuperscript{t} de prés. finie}}} de schémas noethériens. Alors $f$ est un morphisme de descente pour la catégorie des schémas abéliens relatifs.

Si $A$ et $B$ sont deux schémas abéliens sur $S$, on sait en effet que $H=\underline{\operatorname{Hom}}_{\mathrm{gr}}(A,B)$ est représentable, et est net sur $S$, donc il suffit d'appliquer le lemme 1.

\uline{Corollaire} 3 Soit $p$ un nb premier, et soit $f:S'\to S$ comme dans le cor. précédent. Alors $f$ est un morphisme de descente pour la catégorie des \deleted{gr} schémas en groupes $p$-divisibles.

L'argument \deleted{\unsure{qui|pr}} du lemme 1 nous ramène à prouver l'analogue du corollaire 1, quand $X=\underline{\operatorname{Hom}}_{\mathrm{gr}}(A,B)$, $A,B$ deux groupes $p$-divisibles. Il faut prouver que $X$ est formellement net.
% archive sheet number: 24

